% Egyptian Hieroglyphs
\section{Egyptian Hieroglyphic Determinatives}

\subsection{Functional Analogy}

Egyptian hieroglyphs include determinatives: signs that classify or specify preceding phonetic signs without being pronounced themselves.

\subsubsection{Determinative Function}
\begin{itemize}
    \item Semantic classification (animate, inanimate, abstract)
    \item Disambiguation (homophone resolution)
    \item Domain specification (religious, administrative, military)
\end{itemize}

\subsection{Structural Parallels}

\subsubsection{Position}
\begin{itemize}
    \item Determinatives follow phonetic signs
    \item Terminal or near-terminal position
    \item Optional but common in formal writing
\end{itemize}

\subsubsection{Non-phonetic}
\begin{itemize}
    \item Not pronounced in reading
    \item Operational rather than phonetic function
    \item Classificatory role
\end{itemize}

\subsection{Differences from Indus}

Key distinctions:
\begin{itemize}
    \item Egyptian system includes phonetic signs (Indus may not)
    \item Determinatives supplement phonetic core (Indus may be purely operational)
    \item Egyptian is clearly linguistic (Indus is ambiguous)
\end{itemize}

\subsection{Relevance}

This comparison suggests:
\begin{itemize}
    \item Non-phonetic signs can be systematically organized
    \item Operational signs can coexist with (or replace) phonetic signs
    \item Administrative contexts favor operational notation
\end{itemize}
